% Part: first-order-logic
% Chapter: natural-deduction
% Section: quantifier-rules

\documentclass[../../../include/open-logic-section]{subfiles}

\begin{document}

\olfileid{fol}{ntd}{qrl}

\olsection{قواعد المكممات}

\subsection{قواعد~$\lforall$}

\begin{defish}
\AxiomC{$!A({\OLId{latin-ordinary}{a}})$}
\RightLabel{\Intro{\lforall}}
\UnaryInfC{$\lforall[{\OLId{latin-ordinary}{x}}][\Atom{!A}{x}]$}
\DisplayProof
\hfill
\AxiomC{$\lforall[{\OLId{latin-ordinary}{x}}][\Atom{!A}{x}]$}
\RightLabel{\Elim{\lforall}}
\UnaryInfC{$!A({\OLId{latin-ordinary}{t}})$}
\DisplayProof
\end{defish}

في قاعدتي~$\lforall$، الحد~${\OLId{latin-ordinary}{t}}$ مغلق، أي لا متغير فيه؛ و~${\OLId{latin-ordinary}{a}}$ هو
!!a{constant} لا يقع في النتيجة
$\lforall[{\OLId{latin-ordinary}{x}}][!A({\OLId{latin-ordinary}{x}})]$، ولا في فرض حاله !!{undischarged} ضمن
ال!!{derivation} المنتهي بـ~$!A({\OLId{latin-ordinary}{a}})$.
ويسمى~${\OLId{latin-ordinary}{a}}$ \emph{المتغير المميَّز} لاستدلال~\Intro{\lforall}.
\footnote{بقيت تسمية «المتغير المميَّز» لأسباب تاريخية، وإن كان~${\OLId{latin-ordinary}{a}}$ في هذه القاعدة رمزًا ثابتًا.}

\subsection{قواعد~$\lexists$}

\begin{defish}
\AxiomC{$\Atom{!A}{t}$}
\RightLabel{\Intro{\lexists}}
\UnaryInfC{$\lexists[{\OLId{latin-ordinary}{x}}][\Atom{!A}{x}]$}
\DisplayProof
\hfill
\AxiomC{$\lexists[{\OLId{latin-ordinary}{x}}][\Atom{!A}{x}]$}
\AxiomC{[$\Atom{!A}{a}$]$^{\OLId{latin-ordinary}{n}}$}
\DeduceC{$!C$}
\DischargeRule{\Elim{\lexists}}{n}
\BinaryInfC{$!C$}
\DisplayProof
\end{defish}

وكذلك~${\OLId{latin-ordinary}{t}}$ هنا حد مغلق. وأما~${\OLId{latin-ordinary}{a}}$ فهو !!a{constant} شرطه ألا يقع في
$\lexists[{\OLId{latin-ordinary}{x}}][!A({\OLId{latin-ordinary}{x}})]$، وهي المقدمة، ولا في النتيجة~$!C$، ولا في فرض
!!{undischarged} من ال!!{derivation}s المؤدية إلى المقدمتين، سوى فروض~$!A({\OLId{latin-ordinary}{a}})$.
ويسمى~${\OLId{latin-ordinary}{a}}$ \emph{المتغير المميَّز} لاستدلال
\Elim{\lexists}.

هذا المنع لورود المتغير المميَّز في النتائج، أو في فرض
!!{undischarged} من ال!!{derivation}s المؤدية إلى مقدمات
\Intro{\lforall} أو~\Elim{\lexists}، مع استثناء الفروض التي يسقطها
حذف المكمّم الوجودي، يسمى \emph{شرط المتغير المميَّز}.

\textbf{تنبيه تصحيحي على الأصل.} أُظهر في هذا الملخص استثناء الفروض
المؤقتة التي تسقطها قاعدة حذف المكمّم الوجودي؛ فقد نصت عليه القاعدة المفصلة
قبله، وإسقاطه من الملخص يوسّع شرط المتغير المميز خطأً.

\begin{explain}
تذكّر الاصطلاح القائل إننا عندما تكون~$!A$ !!a{formula} يكون فيها
ال!!{variable}~${\OLId{latin-ordinary}{x}}$ حرًا، نشير إلى ذلك بكتابة~$!A({\OLId{latin-ordinary}{x}})$. وفي السياق نفسه،
تكون~$!A({\OLId{latin-ordinary}{t}})$ اختصارًا لـ~$\Subst{!A}{{\OLId{latin-ordinary}{t}}}{{\OLId{latin-ordinary}{x}}}$. ولذلك يمكننا أيضًا كتابة
قاعدة~$\Intro\lexists$ هكذا:
\begin{prooftree}
\AxiomC{$\Subst{!A}{{\OLId{latin-ordinary}{t}}}{{\OLId{latin-ordinary}{x}}}$}
\RightLabel{\Intro{\lexists}}
\UnaryInfC{$\lexists[{\OLId{latin-ordinary}{x}}][!A]$}
\end{prooftree}
لاحظ أن~${\OLId{latin-ordinary}{t}}$ قد يرد أصلًا في~$!A$؛ فمثلًا قد تكون~$!A$ هي
$\Atom{\Obj P}{t,x}$. ومن ثم فإن استنتاج
$\lexists[{\OLId{latin-ordinary}{x}}][\Atom{\Obj P}{t,x}]$ من~$\Atom{\Obj P}{t,t}$ تطبيق
صحيح لـ~$\Intro\lexists$---إذ يجوز لك أن ``تستبدل'' بموضع واحد أو
أكثر، وليس بالضرورة كل المواضع، للحد~${\OLId{latin-ordinary}{t}}$ في المقدمة ال!!{variable}
المقيد~${\OLId{latin-ordinary}{x}}$. غير أن شرطي المتغير المميَّز في~$\Intro\lforall$
و$\Elim\lexists$ يقتضيان ألا يرد ال!!{constant}~${\OLId{latin-ordinary}{a}}$ في~$!A$.
ولذلك لا يمكنك أن تستنتج على نحو صحيح
$\lforall[{\OLId{latin-ordinary}{x}}][\Atom{\Obj P}{a,x}]$ من~$\Atom{\Obj P}{a,a}$ باستعمال
$\Intro\lforall$.
\end{explain}

\begin{explain}
أما~\Intro{\lexists} و\Elim{\lforall} فلا يلزمهما شرط المتغير المميَّز؛ ويجوز اختيار الحد المغلق~${\OLId{latin-ordinary}{t}}$ كيف اتفق.
وأما~\Elim{\lexists} و\Intro{\lforall} فشرطهما ألا يقع ال!!{constant}~${\OLId{latin-ordinary}{a}}$ في النتيجة أو في فرض
!!{undischarged}. وهذا شرط السلامة: فلا
!!{derive}s النظام إلا !!{sentence}s تلزم عن فروضه ال!!{undischarged}.
ولو سقط هذا الشرط لجاز الآتي:
\begin{prooftree}
  \AxiomC{$\lexists[{\OLId{latin-ordinary}{x}}][!A({\OLId{latin-ordinary}{x}})]$}
  \AxiomC{$\Discharge{!A({\OLId{latin-ordinary}{a}})}{{\OLMathNumeral{1}}}$}
  \RightLabel{*\Intro{\lforall}}
  \UnaryInfC{$\lforall[{\OLId{latin-ordinary}{x}}][!A({\OLId{latin-ordinary}{x}})]$}
  \RightLabel{\Elim{\lexists}}
  \BinaryInfC{$\lforall[{\OLId{latin-ordinary}{x}}][!A({\OLId{latin-ordinary}{x}})]$}
\end{prooftree}
لكن~$\lexists[{\OLId{latin-ordinary}{x}}][!A({\OLId{latin-ordinary}{x}})] \Entails/ \lforall[{\OLId{latin-ordinary}{x}}][!A({\OLId{latin-ordinary}{x}})]$.

بما أن قواعد الحذف للمكممات لا تسمح إلا بإحلال حدود مغلقة محل
ال!!{variable}s، فإنه يلزم أن تكون أي !!{formula} يمكن اشتقاقها من
مجموعة من ال!!{sentence}s هي نفسها !!a{sentence}.
\end{explain}


\end{document}
